\documentclass[10pt]{extarticle}
\usepackage[a4paper,margin=1.5cm]{geometry}
\usepackage[T1]{fontenc}
\usepackage{amsmath,amssymb}
\usepackage{graphicx}
\usepackage[export]{adjustbox}
\usepackage{multicol}
\usepackage{enumitem}
\usepackage{booktabs,array,tabularx}
\setlength{\columnsep}{0.7cm}
\setlength{\parindent}{0pt}
\setlength{\parskip}{3pt}
\pagestyle{empty}
\makeatletter
% figure registry (figdefs.tex is generated by the extract step)
\newcommand{\qfigdef}[4]{\@namedef{qf@file@#1}{#2}\@namedef{qf@w@#1}{#3}\@namedef{qf@wide@#1}{#4}}
\newcommand{\qf@check}[1]{\@ifundefined{qf@file@#1}{\PackageError{score2}{Unknown figure id #1}{}}{}}
\newcommand{\qf@img}[4][]{\includegraphics[width=#2\dimexpr\@nameuse{qf@w@#4}pt\relax,max width=#3,max height=0.45\textheight,#1]{\@nameuse{qf@file@#4}}}
% Full-width questions: finalize wraps every question that contains a wide figure in \qfullbegin...\qfullend,
% so the WHOLE question (stem, figures, options) is set across the page in reading order, then the two columns
% resume. The question's head - its number, stem and first figure (\qfullsplit marks the end of the head) -
% is kept on one page, so a number is never left at a page foot with its figure overleaf; the rest may break
% across pages normally, so no big gaps are left. Without a \qfullsplit, the whole question is kept together
% if it is under 45% of a page. Never write these yourself.
\newif\ifqf@full
\newif\ifqf@boxed
\newsavebox{\qf@box}
\newcommand{\qf@open}{\global\qf@fulltrue\par\global\qf@boxedtrue
  \setbox\qf@box\vbox\bgroup\hsize\textwidth\linewidth\textwidth}
\newcommand{\qf@emit}[1]{\par\egroup\global\qf@boxedfalse
  \ifdim\dimexpr\ht\qf@box+\dp\qf@box\relax<#1\textheight\noindent\usebox{\qf@box}\par\else\unvbox\qf@box\fi}
\newcommand{\qfullsplit}{\ifqf@boxed\qf@emit{0.85}\fi}
\newcommand{\qf@close}{\ifqf@boxed\qf@emit{0.45}\fi\global\qf@fullfalse}
\newcommand{\qfullbegin}{\end{multicols}\qf@open}
\newcommand{\qfullend}{\qf@close\begin{multicols}{2}}
\newcommand{\qfullstart}{\qf@open}
\newcommand{\qfulldone}{\qf@close}
% \qfig[scale]{id}: block figure, centred, directly where it is written (size set by the tool).
% Fallback if a wide figure is not inside a full-width question: break out of the columns in place - no floats.
\newcommand{\qfig}[2][1]{\qf@check{#2}%
  \ifqf@full\qf@blk{#1}{#2}\else
  \ifnum\@nameuse{qf@wide@#2}=1 \expandafter\@firstoftwo\else\expandafter\@secondoftwo\fi
  {\end{multicols}\begin{center}\qf@img{#1}{\textwidth}{#2}\end{center}\begin{multicols}{2}}%
  {\qf@blk{#1}{#2}}\fi}
\newcommand{\qf@blk}[2]{\par\smallskip{\centering\qf@img{#1}{\linewidth}{#2}\par}\smallskip}
% \qfigi[scale]{id}: inline figure (inside option lists / tables / minipages); never breaks columns.
\newcommand{\qfigi}[2][1]{\qf@check{#2}\qf@img[valign=c]{#1}{\linewidth}{#2}}
\makeatother
% structure macros - lint reads these to check completeness and figure placement
\newcommand{\papersec}[1]{\par\bigskip{\centering\Large\bfseries PAPER-#1\par}\medskip}
\newcommand{\partsec}[2]{\par\medskip\noindent{\large\bfseries PART-#1: #2}\par\smallskip}
\newcommand{\sect}[1]{\par\smallskip\noindent\textbf{#1}\par}
\newcommand{\qn}[1]{\par\smallskip\noindent\textbf{#1.}~}
\newcommand{\paraq}[2]{\par\smallskip\noindent\textbf{Paragraph for Questions #1 and #2}\par}
\newenvironment{qoptions}{\begin{itemize}[leftmargin=1.8em,itemsep=1pt,topsep=2pt,label={}]}{\end{itemize}}
\newcommand{\opt}[2]{\item[(#1)] #2}
% figure registry (generated by score2_tools.py finalize)
\qfigdef{p1_phy_s1_q03_1}{Score Advanced - Test 13 (13 May 2023)-img/p1_phy_s1_q03_1.png}{124.1}{0}
\qfigdef{p1_phy_s1_q04_1}{Score Advanced - Test 13 (13 May 2023)-img/p1_phy_s1_q04_1.png}{84.6}{0}
\qfigdef{p1_phy_s1_q05_1}{Score Advanced - Test 13 (13 May 2023)-img/p1_phy_s1_q05_1.png}{92.6}{0}
\qfigdef{p1_phy_s1_q06_1}{Score Advanced - Test 13 (13 May 2023)-img/p1_phy_s1_q06_1.png}{113.7}{0}
\qfigdef{p1_phy_s1_q08_1}{Score Advanced - Test 13 (13 May 2023)-img/p1_phy_s1_q08_1.png}{101.3}{0}
\qfigdef{p1_phy_s1_q09_1}{Score Advanced - Test 13 (13 May 2023)-img/p1_phy_s1_q09_1.png}{158.5}{0}
\qfigdef{p1_phy_s1_q10_1}{Score Advanced - Test 13 (13 May 2023)-img/p1_phy_s1_q10_1.png}{178.5}{0}
\qfigdef{p1_phy_s2_q01_1}{Score Advanced - Test 13 (13 May 2023)-img/p1_phy_s2_q01_1.png}{102.7}{0}
\qfigdef{p1_phy_s2_q02_1}{Score Advanced - Test 13 (13 May 2023)-img/p1_phy_s2_q02_1.png}{68.8}{0}
\qfigdef{p1_phy_s2_q03_1}{Score Advanced - Test 13 (13 May 2023)-img/p1_phy_s2_q03_1.png}{131.3}{0}
\qfigdef{p1_phy_s2_q04_1}{Score Advanced - Test 13 (13 May 2023)-img/p1_phy_s2_q04_1.png}{145.4}{0}
\qfigdef{p1_phy_s2_q05_1}{Score Advanced - Test 13 (13 May 2023)-img/p1_phy_s2_q05_1.png}{118.8}{0}
\qfigdef{p1_phy_s2_q06_1}{Score Advanced - Test 13 (13 May 2023)-img/p1_phy_s2_q06_1.png}{96.8}{0}
\qfigdef{p1_phy_s2_q08_1}{Score Advanced - Test 13 (13 May 2023)-img/p1_phy_s2_q08_1.png}{97.9}{0}
\qfigdef{p1_chem_s1_q01_1}{Score Advanced - Test 13 (13 May 2023)-img/p1_chem_s1_q01_1.png}{245.2}{0}
\qfigdef{p1_chem_s2_q01_1}{Score Advanced - Test 13 (13 May 2023)-img/p1_chem_s2_q01_1.png}{119.9}{0}
\qfigdef{p1_chem_s2_q01_2}{Score Advanced - Test 13 (13 May 2023)-img/p1_chem_s2_q01_2.png}{87.2}{0}
\qfigdef{p1_chem_s2_q01_3}{Score Advanced - Test 13 (13 May 2023)-img/p1_chem_s2_q01_3.png}{82.5}{0}
\qfigdef{p1_chem_s2_q01_4}{Score Advanced - Test 13 (13 May 2023)-img/p1_chem_s2_q01_4.png}{110.8}{0}
\qfigdef{p1_chem_s2_q01_6}{Score Advanced - Test 13 (13 May 2023)-img/p1_chem_s2_q01_6.png}{97.2}{0}
\qfigdef{p1_chem_s2_q01_8}{Score Advanced - Test 13 (13 May 2023)-img/p1_chem_s2_q01_8.png}{118.8}{0}
\qfigdef{p1_chem_s2_q02_1}{Score Advanced - Test 13 (13 May 2023)-img/p1_chem_s2_q02_1.png}{175.3}{0}
\qfigdef{p1_chem_s2_q02_2}{Score Advanced - Test 13 (13 May 2023)-img/p1_chem_s2_q02_2.png}{221.8}{0}
\qfigdef{p1_chem_s2_q04_1}{Score Advanced - Test 13 (13 May 2023)-img/p1_chem_s2_q04_1.png}{130.5}{0}
\qfigdef{p2_phy_s1_para1_2_1}{Score Advanced - Test 13 (13 May 2023)-img/p2_phy_s1_para1_2_1.png}{76.8}{0}
\qfigdef{p2_phy_s1_para3_4_1}{Score Advanced - Test 13 (13 May 2023)-img/p2_phy_s1_para3_4_1.png}{213.9}{0}
\qfigdef{p2_phy_s2_q05_1}{Score Advanced - Test 13 (13 May 2023)-img/p2_phy_s2_q05_1.png}{97.9}{0}
\qfigdef{p2_phy_s2_q06_1}{Score Advanced - Test 13 (13 May 2023)-img/p2_phy_s2_q06_1.png}{92.6}{0}
\qfigdef{p2_phy_s2_para7_8_1}{Score Advanced - Test 13 (13 May 2023)-img/p2_phy_s2_para7_8_1.png}{108.3}{0}
\qfigdef{p2_phy_s2_q07_A}{Score Advanced - Test 13 (13 May 2023)-img/p2_phy_s2_q07_A.png}{108.2}{0}
\qfigdef{p2_phy_s2_q07_B}{Score Advanced - Test 13 (13 May 2023)-img/p2_phy_s2_q07_B.png}{108.3}{0}
\qfigdef{p2_phy_s2_q07_C}{Score Advanced - Test 13 (13 May 2023)-img/p2_phy_s2_q07_C.png}{108.2}{0}
\qfigdef{p2_phy_s2_q07_D}{Score Advanced - Test 13 (13 May 2023)-img/p2_phy_s2_q07_D.png}{108.3}{0}
\qfigdef{p2_phy_s2_para9_10_1}{Score Advanced - Test 13 (13 May 2023)-img/p2_phy_s2_para9_10_1.png}{160.9}{0}
\qfigdef{p2_phy_s3_pre_1}{Score Advanced - Test 13 (13 May 2023)-img/p2_phy_s3_pre_1.png}{46.8}{0}
\qfigdef{p2_phy_s3_pre_2}{Score Advanced - Test 13 (13 May 2023)-img/p2_phy_s3_pre_2.png}{45.6}{0}
\qfigdef{p2_phy_s3_pre_3}{Score Advanced - Test 13 (13 May 2023)-img/p2_phy_s3_pre_3.png}{46.3}{0}
\qfigdef{p2_phy_s3_pre_4}{Score Advanced - Test 13 (13 May 2023)-img/p2_phy_s3_pre_4.png}{46.1}{0}
\qfigdef{p2_phy_s3_q12_1}{Score Advanced - Test 13 (13 May 2023)-img/p2_phy_s3_q12_1.png}{133.7}{0}
\qfigdef{p2_phy_s3_q12_2}{Score Advanced - Test 13 (13 May 2023)-img/p2_phy_s3_q12_2.png}{109.7}{0}
\qfigdef{p2_phy_s3_q13_1}{Score Advanced - Test 13 (13 May 2023)-img/p2_phy_s3_q13_1.png}{134.5}{0}
\qfigdef{p2_phy_s4_para3_4_1}{Score Advanced - Test 13 (13 May 2023)-img/p2_phy_s4_para3_4_1.png}{118.8}{0}
\qfigdef{p2_chem_s2_para5_6_1}{Score Advanced - Test 13 (13 May 2023)-img/p2_chem_s2_para5_6_1.png}{320.1}{1}
\qfigdef{p2_chem_s3_para11_1}{Score Advanced - Test 13 (13 May 2023)-img/p2_chem_s3_para11_1.png}{294.1}{1}
\qfigdef{p2_chem_s3_q12_1}{Score Advanced - Test 13 (13 May 2023)-img/p2_chem_s3_q12_1.png}{332.3}{1}
\qfigdef{p2_chem_s3_q13_2}{Score Advanced - Test 13 (13 May 2023)-img/p2_chem_s3_q13_2.png}{197.6}{0}
\qfigdef{p2_chem_s4_para1_2_1}{Score Advanced - Test 13 (13 May 2023)-img/p2_chem_s4_para1_2_1.png}{92.5}{0}
\qfigdef{p2_chem_s4_para3_4_1}{Score Advanced - Test 13 (13 May 2023)-img/p2_chem_s4_para3_4_1.png}{299.5}{1}
\raggedcolumns
\makeatletter\@beginparpenalty=10000\makeatother
\newcolumntype{L}[1]{>{\raggedright\arraybackslash}p{#1}}
\begin{document}
\begin{center}{\Large\bfseries Score Advanced --- Test 13 (13 May 2023)}\end{center}
\begin{multicols}{2}
\raggedcolumns
\papersec{1}
\partsec{1}{PHYSICS}
\sect{SECTION-II}
\qn{1} In 1997, Indian Airforce's most secret supersonic plane `Foxbat Jet', which can fly upto a maximum height of 23000 m and have maximum velocity of 3500 kmph, was flying in Pakistani Airspace at a height of 20,000 m with subsonic velocity. At some point of time, pilot of the plane decided to go supersonic and will fly with a speed of 1020 m/s, this created a sonic boom. Assuming velocity of sound to be 340 m/s. At what distance (in km) from border pilot should have attempted this to safely reach Indian Airspace? Assume he gets noticed when sound reaches the ground.\par
\qn{2} A positronium consists of an electron and a positron revolving about their common centre of mass. The separation between them (in \AA) in its first excited state is :- (Given that radius of H-atom in ground state is 0.53 \AA)\par
\qn{3} In the circuit shown the capacitors are initially uncharged. At t = 0, switch is closed. At time t, the 2 $\mu$F capacitor attains a charge 20 $\mu$C. The heat produced (in $\mu$J) in the 3$\Omega$ resistance during time t is :-\par
\qfig{p1_phy_s1_q03_1}
\qn{4} Two uniform equal stiff rods AB and AC of mass 7 kg each are hinged at A and placed on a smooth horizontal table. The hinge is not fixed to the table and so both the rods can move on the table. A horizontal blow of 35 Ns is delivered perpendicular to AC at C. Find the resulting linear velocity (in m/s) of the center of mass of rod AC, immediately after impact.\par
\qfig{p1_phy_s1_q04_1}
\qn{5} The figure shows a series LCR circuit connected to variable frequency 250V source with L = 50mH, C = 80$\mu$F and R = 12.5$\Omega$. Find Q factor of the circuit for resonance.\par
\qfig{p1_phy_s1_q05_1}
\qn{6} In the following infinite network circles are made from identical uniform wires. The largest innermost circle is made from the wire of total resistance R and circles next to inner ones are half in radii and joined symmetrically. The outer most circles of the infinite arms are maintained at the potentials shown. The potential difference between C and B (in volts) is :-\par
\qfig{p1_phy_s1_q06_1}
\qn{7} N$^{13}$ is a positron emitter with an Q-value of 1.2 MeV. Find the threshold energy (in MeV) of the reaction p + C$^{13}$ $\rightarrow$ N$^{13}$ + n, if the neutron-hydrogen atom mass difference is 0.78 MeV. [m$_\text{e}$ = 0.51 MeV]\par
\qn{8} For the ac circuit shown in the figure, R = 100 k$\Omega$ and C = 10 pF, the phase difference between V$_\text{in}$ and V$_\text{out}$ is 90$^\circ$ at the input signal frequency of \_\_\_\_\_\_\_\_ kHz. (Round off to 2 decimal places)\par
\qfig{p1_phy_s1_q08_1}
\qn{9} In figure, the mass M = 1 kg is rigidly attached to a massless L shaped rod. The spring modulus k = 900 N/m. The angular frequency of the system (in rad/s) is :-\par
\qfig{p1_phy_s1_q09_1}
\qn{10} A rubber ball is projected from a toy cannon as shown. It hits an inclined plane at a distance L from the point of projection. The colliding ball strikes the plane normally. It returns back to a distance $\ell$ from point of collision. The coefficient of restitution e is :- (Take $\ell$ = 1m, L = 4m, $\alpha$ = 37$^\circ$)\par
\qfig{p1_phy_s1_q10_1}
\sect{SECTION-III}
\qn{1} Two different vertical positions (a) \& (b) of a capillary tube are shown in figure with the lower end inside water. For position (a), contact angle is $\dfrac{\pi}{4}$ rad \& water rises to height h above the surface of water while for position (b) height of the tube outside water is kept insufficient \& equal to $\dfrac{h}{\sqrt{2}}$ then contact angle becomes $\dfrac{\pi}{\text{x}}$ rad, then x is :-\par
\qfig{p1_phy_s2_q01_1}
\qn{2} Two identical rectangular rods of metal are welded end to end in series between temperature 0$^\circ$C and 100$^\circ$C and 10 J of heat is conducted (in steady state process) through the rod in 2.00 min. If 5 such rods are taken and joined as shown in figure maintaining the same temperature difference between A and B, then the time (in min) in which 20 J heat will flow through the rods is :-\par
\qfig{p1_phy_s2_q02_1}
\qn{3} The torsion balance shown in figure is used in an experiment to measure radiation pressure. The torque constant (elastic restoring torque) of the suspension fibre is 1 $\times$ 10$^{-11}$ N-m/deg, and the length of the horizontal rod is 6 cm. The beam from a 3 mW He-Ne laser is incident on the black disc and the mirror disc is completely shielded. The angle between the equilibrium positions of the horizontal bar when the beam is switched from ``off'' to ``on'' is $\alpha$ $\times$ 10$^{-2}$ degree. The value of $\alpha$ is :-\par
\qfig{p1_phy_s2_q03_1}
\qn{4} In a water supply system, horizontal nozzles are attached to the main supply pipe at regular interval of 10m. Each nozzle discharges water horizontally at the rate of 2 kg/s with velocity 5 m/s. An empty tank of mass 100 kg and length 20m is placed over smooth horizontal rails below the supply line initially at the rest. Find the velocity of the tank (in m/s) after 100 sec. Ignore vertical gap between the tank and water supply line. Water started moving with the tank as it falls in it.\par
\qfig{p1_phy_s2_q04_1}
\qn{5} A sound wave propagating along x-axis, in medium I of density $\rho_1$ = 1.5 kg/m$^3$ is transmitted to a medium II of density $\rho_2$ = 3 kg/m$^3$ as shown.\par
\qfig{p1_phy_s2_q05_1}
The equation of excess pressure developed by wave in medium I and that in medium II respectively are\par
$p_1 = 4\times10^{-2}\cos\omega\left(\text{t}-\dfrac{\text{x}}{400}\right)$ (In SI unit)\par
$p_2 = 3\times10^{-2}\cos\omega\left(\text{t}-\dfrac{\text{x}}{1200}\right)$ (In SI unit)\par
Then the ratio of intensity of transmitted wave I$_2$ (wave in medium II) to the intensity of incident wave I$_1$ (wave in medium I), that is $\dfrac{\text{I}_2}{\text{I}_1}$ is $\dfrac{\text{n}}{32}$ find n.\par
\qn{6} In the circuit diagram a capacitor which is initially uncharged is connected to an ideal cell of emf E = 5 volt through a resistor R. A leaky dielectric fills the space between the plates of capacitor. The capacitance of capacitor with dielectric is C = 2$\mu$F. Find charge at capacitor in steady state. Resistance of dielectric is also R.\par
\qfig{p1_phy_s2_q06_1}
\qn{7} A large concave mirror with a radius of curvature of 0.5 m is placed on a horizontal plate. We pour a small amount of water into the mirror and place on it a plano-convex lens with a radius of curvature of 1.0 m whose plane side just covers the water. Then we place a point light source above this setup at a height of 0.3 m. The distance (in meters) from the object where the final image is formed will be $\dfrac{\text{x}}{90}$. Find x. (Given that the refractive index of water is $\dfrac{4}{3}$ and the refractive index of glass is 1.50)\par
\qn{8} A uniform electric field E is present horizontally along the paper throughout region but uniform magnetic field B$_0$ is present horizontally (perpendicular to plane of paper in inward direction) right to the line AB as shown. A charge particle having charge q and mass m is projected vertically upward and crosses the line AB after time t$_0$ = $\dfrac{3}{20}$ seconds. Find the speed (in m/s) of projection if particle moves after t$_0$ with constant velocity. (Given : qE = mg)\par
\qfig{p1_phy_s2_q08_1}
\partsec{2}{CHEMISTRY}
\sect{SECTION-II}
\qn{1} Rate law for substitution of 2-bromobutane with OH$^-$ in 75\% ethanol and 25\% water at 30$^\circ$C is given as :\par
\qfig{p1_chem_s1_q01_1}
Calculate \% of reaction taking place by S$_\text{N}$2 mechanism when [OH$^-$] = 1M\par
\qn{2} Electrolysis of acetic acid was carried out in aqueous basic medium. Assuming 75\% yield of reaction, 4.4 gm of CO$_{2(\text{g})}$ was produced at anode. Calculate sum of masses in gms for other gases obtained at anode and cathode.\par
\qn{3} The density of a gas filled in an electric lamp is 0.75 kg/m$^3$. After the lamp has been switched on for sufficient time, the pressure of gas in it increased from 4 $\times$ 10$^4$ Pa to 9 $\times$ 10$^4$ Pa. The increase in u$_\text{rms}$ of gas molecules (in ms$^{-1}$) is.\par
\qn{4} A conductivity cell, when filled with an aqueous solution of 0.02M-KCl at 25$^\circ$C, had a resistance of 240 ohm. Its resistance, when filled with 6 $\times$ 10$^{-5}$ M-NH$_4$OH solution, was 10$^5$ ohm. Calculate the degree of dissociation of NH$_4$OH. Given : Conductivity of 0.02M-KCl solution was 0.2714 S m$^{-1}$.\par
Molar conductances at infinite dilution of NH$_4^+$ and OH$^-$ ions are 7.34 $\times$ 10$^{-3}$ and 1.98 $\times$ 10$^{-2}$ S m$^2$ mol$^{-1}$, respectively.\par
\qn{5} A cloud mass moving across the ocean at an altitude of 2000 m encounters a coastal mountain range. As it rises to a height of 3500 m to pass over the mountains, it undergoes an adiabatic expansion. The pressure at 2000 m and 3500 m is 0.80 and 0.60 atm, respectively. If the initial temperature of the cloud mass is 290 K, what is the cloud temperature (in K) as it passes over the mountains ? Assume that C$_\text{P, m}$ for air is 29.88 JK$^{-1}$ mol$^{-1}$ and that air obeys the ideal gas law. (R = 8.3 J K$^{-1}$ mol$^{-1}$, log 2 = 0.30, log 3 = 0.48, 30 $\times$ log 0.92 = $-$1).\par
\qn{6} At 25$^\circ$C, the vapour pressure of water increases from 30 mm Hg to 31 mm Hg on increasing the external pressure up to 10 atm. If the vapour pressure of water is to be maintained at 30 mm Hg at 25$^\circ$C and 10 atm external pressure, how many gram NaCl should be added per kg of water ?\par
\qn{7} The pH of 0.05 M-Fe(NO$_3$)$_2$ solution is (K$_\text{h}$ of Fe$^{2+}$ = 5.0 $\times$ 10$^{-9}$, log 2 = 0.30).\par
\qn{8} (2-ethylanthraquinol) $\xrightarrow[\text{oxidized in stream of air}]{\text{non catalytically}}$ (X) + $\underset{\text{oxidized product}}{\text{(Y)}}$\par
Find number of hydrogen atoms in oxidized product of above reaction.\par
\qn{9} A lanthanide metal (Ln) have atomic radii of 199 pm and its Ln$^{+3}$ ion have 95 pm which is larger as compared to size of Lu$^{+3}$ \& smaller than size of Ce$^{+3}$ ion. +2 ion of Lanthanide Ln has exactly half filled 4f subshell. Find total number of electrons in its +3 ion.\par
\qn{10} A 10 ml, `x' M solution of Cu(NO$_3$)$_2$ reacts with 40 ml, 0.2 M KCN to produce `y' moles of white ppt and 44.8 ml of a gas at 1 atm and 273 K. Addition of a minimum of 50 ml, `z' M KCN solution was required to completely dissolve the white precipitate. If oxidation state of `C' in the gas evolved is `p', calculate (x + y + z + p)\par
\sect{SECTION-III}
\qn{1} In how many reactions alcohol is obtained as a major product :\par
\qfig{p1_chem_s2_q01_1}
\qfig{p1_chem_s2_q01_2}
\qfig{p1_chem_s2_q01_3}
\qfig{p1_chem_s2_q01_4}
(v) (Me)$_3$CH $\xrightarrow[\Delta]{\text{KMnO}_4/\text{H}^+}$\par
\qfig{p1_chem_s2_q01_6}
(vii) (MeCOO)$_2$Ca $\xrightarrow[\text{distillation}]{\text{dry}}$ $\xrightarrow{\text{NaBH}_4}$\par
\qfig{p1_chem_s2_q01_8}
\qn{2} Consider the following compounds and count number of compounds which can produce tribromoderivatives on reaction with Br$_2$/H$_2$O.\par
\qfig{p1_chem_s2_q02_1}
\qfig{p1_chem_s2_q02_2}
\qn{3} 1 mole of a polypeptide of hydrolysis produced 4 moles of glycine, 2 moles of alanine, 1 mole of phenyl alanine and 1 mole of leucine.\par
If x = number of peptide linkages in polypeptide molecule\par
y = possible number of polypeptide if all identical amino acids remain together in all polypeptides.\par
z = Mass of water (gm) gained by 1 mole of polypeptide on complete hydrolysis.\par
Calculate : $\dfrac{yz}{108x}$\par
\qn{4}\par
\qfig{p1_chem_s2_q04_1}
Unknown mass of compound (P) reacted with different reagents to form gases (A), (B) and (C). Most basic gas released in above process was compltely titrated by 100 ml of 0.2M HCl.\par
If W$_\text{A}$ = mass of gas A (gas), in gm\par
W$_\text{B}$ = mass of gas B (gas), in gm\par
W$_\text{C}$ = mass of gas C (gas), in gm\par
Calculate : (W$_\text{A}$ + W$_\text{B}$ + W$_\text{C}$) $\times$ 10\par
\qn{5} How many of the following statements are correct ?\par
(A) Delta at the mouth of river is formed due to coagulation of the colloidal sea water by the electrolytes present in the river water.\par
(B) Artificial rain is caused by the coagulation of colloidal water droplets in the cloud by an electrolyte.\par
(C) Cottrell precipitator is used in the precipitation of the colloidal waste in the sewage water.\par
(D) Cheese is an example of solid foam.\par
(E) A colloid formed by adding small amount of AgNO$_3$ to KI solution is negatively charged.\par
(F) Micelles are formed only below the Kraft temperature.\par
(G) Mixing of hydrated ferric oxide sol and arsenious sulphide sol bring mutual coagulation.\par
(H) An oil in water emulsion can be diluted by addition of water.\par
(I) Emulsions show Brownian motion and Tyndall effect.\par
\qn{6} Hapticity of cycloheptatriene in Mo(C$_7$H$_8$) (CO)$_3$ is \_\_\_\_\_\_\_\_\par
(Atomic No. of Mo = 42)\par
\qn{7} Identify number of species which have two types of bond length data in their bonding parameters.\par
PCl$_3$ , PCl$_5$ , POCl$_3$ , Cr$_2$O$_7^{2-}$ , MnO$_4^-$ , N$_2$O , NO$_2$ , P$_4$ , C$_{60}$ , BrF$_3$ , BrF$_5$ , IF$_7$ , Al$_2$Cl$_6$\par
\qn{8} Identify number of molecules which have C$_2$-axis of symmetry and they are optically active.\par
cis[Pt(NH$_3$)$_2$Cl$_2$] , 1,3dibromoallene ,\par
cis[Co(en)$_2$Cl$_2$]Cl , trans[Co(en)$_2$Br$_2$]NO$_3$ ,\par
trans 1-2dichlorocyclopentane ,\par
Na$_3$[Fe(C$_2$O$_4$)$_3$] , K[Co(EDTA)]\par
cis[Pt(NH$_3$)$_4$Cl$_2$]Br$_2$\par
\partsec{3}{MATHEMATICS}
\sect{SECTION-II}
\qn{1} Let for x $\in$ $\left(0,\dfrac{\pi}{2}\right)$, f(x) = cot$^{-1}$(2 $-$ tanx) and g($\theta$) = $\displaystyle\int_{\min.\{\theta,\,\pi/4\}}^{\max.\{\theta,\,\pi/4\}}$ f(x) dx then the value of $\left(\dfrac{\pi^2}{g\left(\frac{\pi}{12}\right)-g\left(\frac{\pi}{3}\right)}\right)$ is\par
\qn{2} Given a, b $\in$ \{0, 1, 2, 3, ........, 19, 20\}, consider the system of equations\par
x + y + z = 4\par
2x + y + 3z = 6\par
x + 2y + az = b\par
Let\par
P denotes number of ordered pairs (a, b) so that system of equations has a unique solution\par
Q denotes number of ordered pairs (a, b) so that system of equations has no solution\par
R denotes number of ordered pairs (a, b) so that system of equations has infinite solutions,\par
then (P + Q + R) is :\par
\qn{3} If $|2z_1+\bar{z}_2| = 2\sqrt{2}$ \& $|1+2z_1z_2| = 3$ then minimum value of $16(|z_1|^2+4|z_2|^2)$, is\par
\qn{4} Let a = $\displaystyle\sum_{r=1}^{11}\sin^2\left(\dfrac{r\pi}{24}\right)$ ; b = $\displaystyle\sum_{r=1}^{11}\cos^2\left(\dfrac{r\pi}{24}\right)$ ; c = $\displaystyle\sum_{r=1}^{11}\tan^3\left(\dfrac{r\pi}{24}\right)$ and d = $\displaystyle\sum_{r=1}^{11}\cot^3\left(\dfrac{r\pi}{24}\right)$\par
Then the value of a$^3$ + b$^3$ $-$ 11 (a + b) + 33ab + log$_{(2d-c)}$(2c$-$d) is\par
\qn{5} Let image of the line $\dfrac{x-1}{3}=\dfrac{y-3}{5}=\dfrac{z-4}{2}$ in the plane 2x $-$ y + z + 3 = 0 be L.\par
A plane 7x + py + qz + r = 0 (p, q, r $\in$ R) is such that it contains the line L and perpendicular to the plane 2x $-$ y + z + 3 = 0. Then the value of (p + q + r) is equal to\par
\qn{6} The value of $\displaystyle\sum_{r=0}^{10}(-1)^r4^{10-r}\,{}^{30}\mathrm{C}_r\cdot{}^{30-r}\mathrm{C}_{10-r}$ is equal to ${}^{30}\mathrm{C}_k\cdot 3^k$ and $\displaystyle\sum_{r=0}^{20}{}^{30}\mathrm{C}_r\,{}^{20}\mathrm{C}_r={}^{n}\mathrm{C}_{20}$ then the value of (n + k) is\par
\qn{7} If the probability that a number selected from the set \{1, 2, 3, ..., 1000\} is divisible by 3 but neither divisible by 5 nor by 7, is $\dfrac{m}{n}$, then (n $-$ m) is (where m and n are relatively prime numbers)\par
\qn{8} Let S = $\dfrac{1^2}{1.3}+\dfrac{2^2}{3.5}+\dfrac{3^2}{5.7}+\ldots+\dfrac{500^2}{999.1001}$, then [S] is\par
(where [.] denotes greatest integer function)\par
\qn{9} The function y = $f$(x) is the solution of the differential equation $\dfrac{dy}{dx}+\dfrac{xy}{x^2-4}=\dfrac{7x^6+2x}{\sqrt{4-x^2}}$ in ($-$2,2) satisfying $f$(0) = 1. If $\displaystyle\int_{-\sqrt{3}}^{\sqrt{3}} f(x)\,dx = a\pi+b\sqrt{3}$, a,b are rational numbers, then the value of (6a + b) is\par
\qn{10} Let x + y = 0 and x $-$ y = 0 are tangent to a parabola whose focus is S(1, 2).AB be a focal chord of the parabola.If the harmonic mean of AS and BS can be expressed as $\dfrac{m}{\sqrt{n}}$ (where m and n are coprime) then the value of (m + 3n) is equal to\par
\sect{SECTION-III}
\qn{1} Let f be a twice differentiable function with f($-$1) = 1 = f(1) and f(0) = 0. Then the minimum number of root(s) of equation\par
\resizebox{\linewidth}{!}{$\mathrm{f(x)(f\,'(x))^3+(f(x))^2\,f\,'(x)\,f\,''(x)-f(x)\,f\,'(x)=x(f\,'(x))^2+x\,f(x)\,f\,''(x)-x}$}\par
is(are)\par
\qn{2} Let $\alpha$ be a real root of the equation, x$^5$ $-$ x$^3$ + x $-$ 2 = 0. Then the value of [$\alpha^6$] is (where [*] denotes the greatest integer function)\par
\qn{3} The volume of parallelopiped whose coterminous edges are represented by the vectors $3\left(\vec{b}\times\vec{c}\right)$, $2\left(\vec{a}\times\vec{b}\right)$ and $4\left(\vec{c}\times\vec{a}\right)$ where\par
$\vec{a}=(1+\sin\theta)\hat{i}+(\cos\theta)\hat{j}+(\sin2\theta)\hat{k}$\par
$\vec{b}=\sin\left(\theta+\dfrac{2\pi}{3}\right)\hat{i}+\cos\left(\theta+\dfrac{2\pi}{3}\right)\hat{j}+\sin\left(2\theta+\dfrac{4\pi}{3}\right)\hat{k}$\par
$\vec{c}=\sin\left(\theta-\dfrac{2\pi}{3}\right)\hat{i}+\cos\left(\theta-\dfrac{2\pi}{3}\right)\hat{j}+\sin\left(2\theta-\dfrac{4\pi}{3}\right)\hat{k}$\par
is 18 cubic units. Then the number of value(s) of $\theta\in\left(0,\dfrac{\pi}{2}\right)$ is(are)\par
\qn{4} A continuous function f : R $\rightarrow$ R satisfies the functional equation f(4x) $-$ 2f(2x) + f(x) = x$^2$ such that f(0) = 1 Then the value of f($3\sqrt{3}$) is\par
\qn{5} Let AB be a focal chord of an ellipse $\dfrac{x^2}{16}+\dfrac{y^2}{9}=1$ passing through one focus S and S$'$ be its other focus, then the value of semi-perimeter of $\Delta$S$'$AB is\par
\qn{6} If\par
\resizebox{\linewidth}{!}{$\displaystyle\int e^{e^x+x^3+x}\left[\left(xe^x+3x^3+x+1\right)\sin x+x\cos x\right]dx = f(x)e^{e^x+x^3+x}+C$}\par
(where $f\left(\dfrac{\pi}{6}\right)=\dfrac{\pi}{12}$ and C is the constant of integration and [*] denotes Greatest Integer Function) then $\left[f\left(\dfrac{5\pi}{2}\right)\right]$ is equal to\par
\qn{7} Number of solution(s) of the equation 8000 sin$^3$x + 3375 cos$^3$ x + 1728 = 5400 sin 2x in (0, 2$\pi$) is(are)\par
\qn{8} If $\displaystyle\lim_{x\to2}\dfrac{2+x+x^2+x^3+\ldots\ldots\ldots\ldots+x^{n-1}-2^n}{x-2}$ = 4097, then the value of [ $\sqrt{n}$ ] equals (where [*] denotes the greatest integer function)\par
\raggedcolumns
\papersec{2}
\partsec{1}{PHYSICS}
\sect{SECTION-I (i)}
\paraq{1}{2}Consider three large parallel metallic plates out of which middle is given charge Q and other two are neutral (see figure). The outer plates are connected with each other. Space between right plate and middle plate is filled by a medium of dielectric constant K = 3.\par
\qfig{p2_phy_s1_para1_2_1}
\qn{1} The charge appearing on the :-
\begin{qoptions}
\opt{A}{outer surface of left plate is $\dfrac{\text{Q}}{2}$.}
\opt{B}{inner surface of left plate is $-\dfrac{\text{Q}}{2}$.}
\opt{C}{left surface of middle plate is $\dfrac{\text{Q}}{2}$.}
\opt{D}{net charge on left plate is zero.}
\end{qoptions}
\qn{2} Select correct alternative (s) :-
\begin{qoptions}
\opt{A}{Electric field at point A is three times of electric field at point B.}
\opt{B}{Electric field at point A is equal to electric field at point B.}
\opt{C}{Energy stored between left and middle plate is 3 times of energy between right and middle plate.}
\opt{D}{Energy stored between left and middle plate is equal to energy between right and middle plate.}
\end{qoptions}
\paraq{3}{4}A man inside the lift swings a stone tied to a string of length L in a vertical plane. The string remains taut throughout the time of rotation of the stone.\par
\textbf{Case-I :} The lift accelerates vertically upwards with acceleration a$_1$ = $\dfrac{\text{g}}{2}$.\par
\textbf{Case-II :} The lift accelerates horizontally towards right with acceleration a$_2$ = g.\par
\qfig{p2_phy_s1_para3_4_1}
\qn{3} For case-I choose the correct option(s).
\begin{qoptions}
\opt{A}{The tension at the lowest point of its path for just completing the circular motion is 9 mg.}
\opt{B}{The minimum speed `u$_1$' of the stone (relative to the lift) required at the lowest point of its path for successful completion of the loop is $\sqrt{\dfrac{15}{2}\text{gL}}$.}
\opt{C}{The minimum speed `u$_1$' of the stone (relative to the lift) required at the lowest point of its path for successful completion of the loop is $\sqrt{5\text{gL}}$.}
\opt{D}{The tension at the lowest point of its path for just completing the circular motion is 6mg.}
\end{qoptions}
\qn{4} For case-II choose the correct option(s).
\begin{qoptions}
\opt{A}{The maximum possible tension in the string, if the particle just completes the circular motion is 6mg.}
\opt{B}{The minimum speed `u$_2$' of the stone (relative to the lift) required at the lowest point of its path for successful completion of the loop is $\sqrt{(2+3\sqrt{2})\text{gL}}$.}
\opt{C}{The minimum speed `u$_2$' of the stone (relative to the lift) required at the lowest point of its path for successful completion of the loop is $\sqrt{(2+4\sqrt{2})\text{gL}}$.}
\opt{D}{The maximum possible tension in the string, if the particle just completes the circular motion is $6\sqrt{2}$ mg.}
\end{qoptions}
\sect{SECTION-I (ii)}
\paraq{5}{6}In August 2022, world highest rail bridge on river Chenab was completed and inaugurated. It is an arch bridge at a height of 1178 ft above water and arch span of 1532 ft. Arch bridge have been built from a very long time ago. The stability of arch bridge depend on the over hang of bricks used to make to bridge.\par
\qn{5} Four bricks are to be stacked at the edge of a table, each bricks over hanging on the one below it so that the top brick extends as far as possible beyond the edge of the table. The arrangement will be appear like as shown. The value of x$_1$, x$_2$, x$_3$ and x$_4$ will be :-\par
\qfig{p2_phy_s2_q05_1}
\begin{qoptions}
\opt{A}{$\dfrac{\ell}{2},\dfrac{\ell}{2},\dfrac{\ell}{2},\dfrac{\ell}{2}$}
\opt{B}{$\dfrac{\ell}{2},\dfrac{\ell}{4},\dfrac{\ell}{8},\dfrac{\ell}{16}$}
\opt{C}{$\dfrac{\ell}{16},\dfrac{\ell}{8},\dfrac{\ell}{4},\dfrac{\ell}{2}$}
\opt{D}{$\dfrac{\ell}{8},\dfrac{\ell}{6},\dfrac{\ell}{4},\dfrac{\ell}{2}$}
\end{qoptions}
\qn{6} The minimum number of bricks each having length 0.3m required to make an arch bridge of span 1m as shown in figure is :-\par
\qfig{p2_phy_s2_q06_1}
\begin{qoptions}
\opt{A}{15}
\opt{B}{33}
\opt{C}{35}
\opt{D}{37}
\end{qoptions}
\paraq{7}{8}A rectangular wire frame is kept on the principal axis of a concave mirror as shown in figure. F \& C are the focus \& centre of curvature of mirror respectively.\par
\qfig{p2_phy_s2_para7_8_1}
f is the focal length of mirror. A small particle is moving along the length of wire frame with constant velocity. Based on above information answer the following questions.\par
\qn{7} The image of the given object is best shown by :-
\begin{qoptions}
\opt{A}{\qfigi{p2_phy_s2_q07_A}}
\opt{B}{\qfigi{p2_phy_s2_q07_B}}
\opt{C}{\qfigi{p2_phy_s2_q07_C}}
\opt{D}{\qfigi{p2_phy_s2_q07_D}}
\end{qoptions}
\qn{8} Mark the correct statement(s) :-
\begin{qoptions}
\opt{A}{The velocity of image of particle is remaining constant.}
\opt{B}{The velocity of image of particle is varying.}
\opt{C}{The total length of image of wire frame is $\left(\dfrac{7+\sqrt{17}}{4}\right)$f.}
\opt{D}{Both (B) and (C) are correct.}
\end{qoptions}
\paraq{9}{10}Two vertical cylinders of different areas of cross-sections (A$_2$ $<$ A$_1$) shown in the figures have pistons of same mass (m$_1$ = m$_2$) fitted on the top. The cylinders enclose N moles of monatomic gas. The initial pressure and initial volume of the gas are P and V, respectively. It is given that spring I is relaxed. Now through the electric heater, heat is slowly supplied. Finally spring II becomes relaxed. It is given that the cylinders and the pistons are thermally insulated.\par
\qfig{p2_phy_s2_para9_10_1}
\qn{9} What is the final pressure?
\begin{qoptions}
\opt{A}{$\text{P}+\text{mg}\left(\dfrac{\text{A}_1+\text{A}_2}{\text{A}_1\text{A}_2}\right)$}
\opt{B}{$\text{P}+\text{mg}\left(\dfrac{\text{A}_1-\text{A}_2}{\text{A}_1\text{A}_2}\right)$}
\opt{C}{$\text{P}+2\text{mg}\left(\dfrac{\text{A}_1-\text{A}_2}{\text{A}_1\text{A}_2}\right)$}
\opt{D}{$\text{P}-2\text{mg}\left(\dfrac{\text{A}_1+\text{A}_2}{\text{A}_1\text{A}_2}\right)$}
\end{qoptions}
\qn{10} What is the final volume?
\begin{qoptions}
\opt{A}{$\text{V}+\dfrac{\text{mg}}{\text{k}}\dfrac{(\text{A}_1+\text{A}_2)(\text{A}_1^2+\text{A}_2^2)}{\text{A}_1\text{A}_2}$}
\opt{B}{$\text{V}+\dfrac{\text{mg}}{\text{k}}\dfrac{(\text{A}_1-\text{A}_2)(\text{A}_1^2+\text{A}_2^2)}{\text{A}_1\text{A}_2}$}
\opt{C}{$\text{V}+\dfrac{\text{mg}}{2\text{k}}\dfrac{(\text{A}_1-\text{A}_2)(\text{A}_1^2+\text{A}_2^2)}{\text{A}_1\text{A}_2}$}
\opt{D}{$\text{V}+\dfrac{2\text{mg}}{\text{k}}\dfrac{(\text{A}_1-\text{A}_2)(\text{A}_1^2+\text{A}_2^2)}{\text{A}_1\text{A}_2}$}
\end{qoptions}
\sect{SECTION-I (iii)}
\qn{11} \textit{(Paragraph for Question 11)} In List-II some arrangements of point charges with shaded Gaussian surface are given. In List-I the electric flux passing through surface are given. Match the following columns with flux passes through shaded region.\par
{\footnotesize\setlength{\tabcolsep}{1.5pt}\noindent\begin{tabular}{|c|c|c|L{3.3cm}|}
\hline
 & List-I & & List-II \\
\hline
(a) & $\left(\dfrac{\text{Q}}{2\varepsilon_0}\right)$ & (P) & \qfigi[0.75]{p2_phy_s3_pre_1} \\
\hline
(b) & $\left(\dfrac{\text{Q}}{4\varepsilon_0}\right)$ & (Q) & \qfigi[0.75]{p2_phy_s3_pre_2} \\
\hline
(c) & $\left(\dfrac{\text{Q}}{24\varepsilon_0}\right)$ & (R) & \qfigi[0.75]{p2_phy_s3_pre_3} \\
\hline
(d) & $\left(\dfrac{\text{Q}}{48\varepsilon_0}\right)$ & (S) & \qfigi[0.75]{p2_phy_s3_pre_4} \\
\hline
\end{tabular}}\par
\begin{qoptions}
\opt{A}{a$\rightarrow$Q, b$\rightarrow$R, c$\rightarrow$P, d$\rightarrow$S}
\opt{B}{a$\rightarrow$Q, b$\rightarrow$P, c$\rightarrow$R, d$\rightarrow$S}
\opt{C}{a$\rightarrow$P, b$\rightarrow$Q, c$\rightarrow$S, d$\rightarrow$R}
\opt{D}{a$\rightarrow$P, b$\rightarrow$Q, c$\rightarrow$R, d$\rightarrow$S}
\end{qoptions}
\qn{12} \textit{(Paragraph for Question 12)} There are four situations in which is a conducting loop of side L can be kept. In each of the situations, one of the vertices is at origin. If magnitude field in space is $\vec{\text{B}}$ = $\beta_0$[5yzt$\hat{\text{i}}$ + zxt$\hat{\text{j}}$ + 3y$^2$t$\hat{\text{k}}$].\par
\qfig{p2_phy_s3_q12_1}
\begin{qoptions}
\opt{A}{a$\rightarrow$P, b$\rightarrow$Q, c$\rightarrow$R, d$\rightarrow$S}
\opt{B}{a$\rightarrow$P, b$\rightarrow$Q, c$\rightarrow$R, d$\rightarrow$R}
\opt{C}{a$\rightarrow$P, b$\rightarrow$R, c$\rightarrow$Q, d$\rightarrow$S}
\opt{D}{a$\rightarrow$P, b$\rightarrow$Q, c$\rightarrow$S, d$\rightarrow$R}
\end{qoptions}
\qn{13} \textit{(Paragraph for Question 13)} Two identical rods of length $\ell$ having variable linear mass density $\lambda$ = 15x$^2$ (in S.I. units) are joint at an angle 74$^\circ$. Here x is the distance from the joint of the rod. The rods are symmetrically placed over a disc of radius R = 1.8 m in vertical plane as shown in the figure. Consider all surfaces to be smooth. For some values of $\ell$ the rods oscillates in SHM for small angular amplitude in its plane. The time period of oscillation in this case is given by $2\pi\sqrt{\dfrac{\text{y}}{\text{x}}}$, where x and y are coprime. (g = 10 m/s$^2$)\par
\qfig{p2_phy_s3_q12_2}
Match the list for different values of $\ell$.\par
{\footnotesize\setlength{\tabcolsep}{1.5pt}\noindent\begin{tabular}{|c|c|c|L{4.7cm}|}
\hline
 & List-I & & List-II \\
\hline
(a) & $\ell$ = 2m & (P) & System is in unstable equilibrium. \\
\hline
(b) & $\ell$ = 4m & (Q) & Time period of oscillation is T = $2\pi\sqrt{\dfrac{11}{10}}$ sec. \\
\hline
(c) & $\ell$ = 8m & (R) & Time period of oscillation is T = $2\pi\sqrt{\dfrac{31}{30}}$ sec. \\
\hline
(d) & $\ell$ = 10m & (S) & Time period of oscillation is T = $2\pi\sqrt{\dfrac{51}{50}}$ sec. \\
\hline
\end{tabular}}\par
\begin{qoptions}
\opt{A}{a$\rightarrow$P, b$\rightarrow$P, c$\rightarrow$Q, d$\rightarrow$S}
\opt{B}{a$\rightarrow$P, b$\rightarrow$Q, c$\rightarrow$R, d$\rightarrow$S}
\opt{C}{a$\rightarrow$P, b$\rightarrow$R, c$\rightarrow$Q, d$\rightarrow$S}
\opt{D}{a$\rightarrow$P, b$\rightarrow$P, c$\rightarrow$R, d$\rightarrow$Q}
\end{qoptions}
\qn{14} \textit{(Paragraph for Question 14)} A square platform of side length 8m is situated in x-z plane such that it is at 16m from the x-axis and 8m from the z-axis as shown in figure. A particle is projected with velocity $\vec{\text{v}}$ = v$_2\hat{\text{i}}$ + 25$\hat{\text{j}}$ m/s relative to wind from origin and at the same instant the platform starts with acceleration $\vec{\text{a}}$ = (2$\hat{\text{i}}$ + 2.5$\hat{\text{j}}$) m/s$^2$. Wind is blowing with velocity v$_1\hat{\text{k}}$. (g = 10 m/s$^2$)\par
\qfig{p2_phy_s3_q13_1}
{\footnotesize\setlength{\tabcolsep}{1.5pt}\noindent\begin{tabular}{|c|L{4.1cm}|c|c|}
\hline
 & List-I & & List-II \\
\hline
(a) & Least possible values of v$_2$ (in m/s) so that particle hits the platform or edge of platform is & (P) & 4 \\
\hline
(b) & Least possible values of v$_1$ (in m/s) so that particle hits the platform or edge of platform is & (Q) & 6 \\
\hline
(c) & If t is the time (in second) after particle hits the platform then 2t is equal to & (R) & 8 \\
\hline
(d) & Value of displacement with respect to ground (in m) of the particle in y-direction, when v$_2$ has its minimum possible value is (till particle hits platform or edge of platform) & (S) & 20 \\
\hline
\end{tabular}}\par
\begin{qoptions}
\opt{A}{a$\rightarrow$Q, b$\rightarrow$S, c$\rightarrow$R, d$\rightarrow$P}
\opt{B}{a$\rightarrow$Q, b$\rightarrow$P, c$\rightarrow$R, d$\rightarrow$S}
\opt{C}{a$\rightarrow$Q, b$\rightarrow$R, c$\rightarrow$S, d$\rightarrow$P}
\opt{D}{a$\rightarrow$Q, b$\rightarrow$P, c$\rightarrow$S, d$\rightarrow$R}
\end{qoptions}
\sect{SECTION-II}
\paraq{1}{2}A dish antenna with diameter 20m receives (at normal incidence) a radio signal from a distance source. The radio signal is a continuous sinusoidal wave with amplitude E$_0$ = 0.2 $\mu$V/m. Assume that the antenna is perfectly absorbing.\par
\qn{1} The power received by the antenna is x $\times$ 10$^{-14}$ W. The value of x is :-\par
\qn{2} The force exerted by the radio waves on the antenna is x $\times$ 10$^{-22}$ N. The value of x is :-\par
\paraq{3}{4}A railgun is an electrically powered electromagnetic projectile launcher based on similar principles to the homopolar motor. It comprises a pair of parallel conducting rails, along which a sliding armature is accelerated by the electromagnetic effects of a current that flows down one rail, into the armature and then back along the other rail.\par
Now two parallel metal rails are connected to an electrical power supply. When a conductive projectile is inserted between the rails (at the end connected to the power supply), it completes the circuit. Electrons flow from the negative terminal of the power supply up the negative rail, across the projectile, and down the positive rail, back to the power supply.\par
This current makes the railgun behave as an electromagnet, creating a magnetic field inside the loop formed by the length of the rails up to the position of the armature. In accordance with the right-hand rule, the magnetic field circulates around each conductor. Since the current is in the opposite direction along each rail, the net magnetic field between the rails (B) is directed at right angles to the plane formed by the central axes of the rails and the armature. In combination with the current (I) in the armature, this produces a Lorentz force which accelerates the projectile along the rails, away from the power supply.\par
If radius of the rails is 1 mm (which are assumed to be circular in cross section) , the distance between the centre points of the rails is 1 m and the constant current through the system is 10A. ($\ell$n 10 = $\dfrac{7}{3}$)\par
\qfig{p2_phy_s4_para3_4_1}
\qn{3} Velocity (in m/s) of armature having mass 10 mg after 5 sec is :-\par
\qn{4} If all the heat loss during 5s is only through the armature and all the heat is used to increase the temperature of the armature then calculate the change in temperature of armature (in K). (Specific heat is 25 kJ/kg/K, resistance of armature = 0.1 $\Omega$)\par
\paraq{5}{6}A Carnot heat engine draws heat from a reservoir at temperature T$_1$ and rejects heat to another reservoir at temperature T$_3$. The Carnot forward cycle engine drives a Carnot reversed cycle engine or Carnot refrigerator which absorbs heat from reservoir at temperature T$_2$ and rejects heat to a reservoir at temperature T$_3$. The high temperature T$_1$ is 600 K and low temperature T$_2$ is 300 K.\par
\qn{5} If heat supplied to engine Q$_1$ is equal to heat absorbed by the refrigerator Q$_2$, the value of T$_3$ (in K) is :-\par
\qn{6} The COP of carnot refrigerator is :-\par
\raggedcolumns
\partsec{2}{CHEMISTRY}
\sect{SECTION-I (i)}
\paraq{1}{2}A secondary cell after use can be recharged by passing current through it in the opposite direction so that it can be used again. A good secondary cell can undergo a large number of discharging and charging cycles. The most important secondary cell is the lead storage battery commonly used in automobiles and invertors. It consists of a lead anode (negative terminal) and a grid of lead packed with lead dioxide (PbO$_2$) as cathode (positive terminal). A 38\% solution of sulphuric acid is used as an electrolyte.\par
The cell reactions when the battery is in use are given below :\par
\textbf{Anode :} Pb(s) + SO$_4^{2-}$(aq) $\rightarrow$ PbSO$_4$(s) + 2e$^-$\par
\textbf{Cathode :} PbO$_2$(s) + SO$_4^{2-}$(aq) + 4H$^+$(aq) + 2e$^-$ $\rightarrow$ PbSO$_4$(s) + 2H$_2$O($\ell$)\par
i.e., overall cell reaction consisting of cathode and anode reactions is :\par
Pb(s) + PbO$_2$(s) + 2H$_2$SO$_4$(aq) $\rightarrow$ 2PbSO$_4$(s) + 2H$_2$O($\ell$)\par
On charging the battery the reaction is reversed and PbSO$_4$(s) on anode and cathode is converted into Pb and PbO$_2$, respectively.\par
\qn{1} Which of the following fact(s) make the cell rechargeable ?
\begin{qoptions}
\opt{A}{PbSO$_4$(s) predominantly deposits on the electrodes and not in other parts of the cell.}
\opt{B}{PbSO$_4$(s) deposits on both the electrodes.}
\opt{C}{Resistivity of PbSO$_4$(s) deposits (3 $\times$ 10$^9$ $\Omega$ cm) is much higher than that of PbO$_2$(s) (2 $\times$ 10$^{-4}$ $\Omega$ cm).}
\opt{D}{Resistivity of PbO$_2$(s) (2 $\times$ 10$^{-4}$ $\Omega$ cm) is close to that of Pb(s) (0.2 $\times$ 10$^{-4}$ $\Omega$ cm).}
\end{qoptions}
\qn{2} Electrolysis of water may become a competing process during charging of the cell. For the operation of the cell at 298 K, the correct statement(s) is/are :
\begin{qoptions}
\opt{A}{During charging of this cell, positive terminal of the DC power source is connected to the electrode that gets reduced, and negative terminal to the electrode that gets oxidised.}
\opt{B}{The EMF of the cell depends upon the amounts of Pb species in the cell.}
\opt{C}{After a long time of the lead acid cell overcharging, the evolution of oxygen on one terminal and that of hydrogen on the other terminal may result in the formation of an explosive mixture of gases.}
\opt{D}{The total energy stored in the cell is limited by the moles of Pb(s) in the negative electrode or the moles of PbO$_2$(s) in the positive electrode, whichever amount is less.}
\end{qoptions}
\paraq{3}{4}Sodium salt of Pseudo halide X, Y and Z form colourless solution in water. Solution of X decolorizes I$_3^-$ solution in the presence of Y and gives brisk effervescence. Solution of Y gives an intense red colour compound ``R'' on reaction with Fe$^{+3}$ solution. Solution of Z gives intense blue colour ``B'' on reaction with a solution containing Fe$^{+2}$ \& Fe$^{+3}$ ions.\par
\qn{3} Select correct for X, Y and Z?
\begin{qoptions}
\opt{A}{anions in X, Y \& Z are nitrogen containing}
\opt{B}{anions in X, Y \& Z all have two $\sigma$ bonds}
\opt{C}{reaction of ``X'' and I$_3^-$ is a redox change}
\opt{D}{Colourless solution is obtained by reaction of I$_3^-$ \& X, produces violet blue solution with starch}
\end{qoptions}
\qn{4} Select correct for compound ``R'' and ``B''.
\begin{qoptions}
\opt{A}{``B'' is NaFe$^\text{III}$[Fe$^\text{II}$(CN)$_6$]}
\opt{B}{``R'' is [Fe(SCN)]$^{+2}$}
\opt{C}{``R'' is decolorized with excess KF$_\text{aq.}$}
\opt{D}{``R'' have $\sqrt{35}$ B.M. spin magnetic moment}
\end{qoptions}
\qfullbegin
\sect{SECTION-I (ii)}
\paraq{5}{6}\par
\qfig{p2_chem_s2_para5_6_1}
\qfullend
\qn{5} True statement about (E) is?
\begin{qoptions}
\opt{A}{It gives Tollen's test}
\opt{B}{It gives effervescence with NaHCO$_3$}
\opt{C}{It gives Iodoform test}
\opt{D}{It gives CAN test}
\end{qoptions}
\qn{6} Correct statement is
\begin{qoptions}
\opt{A}{When (D) is treated with Na/NH$_3$($\ell$) trans alkene is formed}
\opt{B}{Compound (G) is formed by Cannizaro reaction}
\opt{C}{Compound (G) is formed by aldol condensation}
\opt{D}{(E) \& (G) are identical}
\end{qoptions}
\paraq{7}{8}In water, HF behaves as a weak acid. The reactions involved are :\par
\textbf{Reaction-I :} HF(aq) + H$_2$O($\ell$) $\rightleftharpoons$ H$_3$O$^+$(aq) + F$^-$(aq); K$_1$ = 1.1 $\times$ 10$^{-3}$\par
\textbf{Reaction-II :} HF(aq) + F$^-$(aq) $\rightleftharpoons$ HF$_2^-$(aq); K$_2$ = 0.26\par
Two chemists wanted to determine the acidity constant, K$_1$, using the same HF solution and measuring the pH of solution. One chemist was having the idea of reaction-I only while the other chemist was having the idea of both reactions. Surprisingly, both the chemist determined the same value of K$_1$ (as given above).\par
\qn{7} Which of the following must be incorrect ?
\begin{qoptions}
\opt{A}{[HF]$_0$ = $\sqrt{2\text{K}_1\text{K}_2}$}
\opt{B}{K$_1$ = [H$_3$O$^+$]}
\opt{C}{[HF]$_0$ = 2K$_1$}
\opt{D}{[HF]$_0$ = 2[H$_3$O$^+$]}
\end{qoptions}
\qn{8} The initial concentration of HF solution used is :
\begin{qoptions}
\opt{A}{0.22 M}
\opt{B}{2.2 $\times$ 10$^{-3}$ M}
\opt{C}{3.3 $\times$ 10$^{-3}$ M}
\opt{D}{2.0 M}
\end{qoptions}
\paraq{9}{10}Metal (M) is extracted from an ore (O) which involve separation of non magnetic ore from magnetic gangue which consist of major Impurity (I) by magnetic separation method.\par
Concentrated ore is some times roasted to remove impurities such As \& S as volatile oxides. The roasted ore is heated with coal in furnace at 1500 K to produce molten metal.\par
The crude metal obtained by this method is contaminated by Fe \& other metals. It is therefore remelted on inclined surface. Metal (M) which has lower melting point than other metallic impurities present, melts and flows down leaving behind the less fusible metals.\par
Metal (M) is a soft silvery white metal it is ductile \& can be rolled into foils.\par
\qn{9} Identify ore (O) , process of reduction of ore \& process of refining of metal respectively :-
\begin{qoptions}
\opt{A}{Galena, self reduction, distillation}
\opt{B}{Cassiterite, carbon reduction, liquation}
\opt{C}{Heamatite , carbon reduction , magnetic separation}
\opt{D}{Cassiterite, self reduction , liquation}
\end{qoptions}
\qn{10} Which oxidation number is not observed for any atom in impurity (I).
\begin{qoptions}
\opt{A}{$-$2}
\opt{B}{+2}
\opt{C}{+6}
\opt{D}{+3}
\end{qoptions}
\qfullbegin
\sect{SECTION-I (iii)}
\qn{11} \textit{(Paragraph for Question 11)}\par
\qfig{p2_chem_s3_para11_1}
\qfullsplit
$\Rightarrow$ For the above given scheme \%yield are mentioned along with the arrow.\par
$\Rightarrow$ x, y and z are mass (gm) of CHI$_3$, (R) and (U) respectively.\par
$\Rightarrow$ Match the list and select correct code from given options.\par
(I = 127, Br = 80)\par
{\footnotesize\setlength{\tabcolsep}{1.5pt}\noindent\begin{tabular}{|c|c|c|c|}
\hline
 & List-I & & List-II \\
\hline
(a) & x is & (P) & 4 \\
\hline
(b) & y is & (Q) & 78.8 \\
\hline
(c) & z is & (R) & 5.6 \\
\hline
(d) & n is & (S) & 7.92 \\
\hline
\end{tabular}}\par
Select correct code :\par
\begin{qoptions}
\opt{A}{a$\rightarrow$P, b$\rightarrow$Q, c$\rightarrow$R, d$\rightarrow$S}
\opt{B}{a$\rightarrow$Q, b$\rightarrow$R, c$\rightarrow$S, d$\rightarrow$P}
\opt{C}{a$\rightarrow$Q, b$\rightarrow$S, c$\rightarrow$P, d$\rightarrow$P}
\opt{D}{a$\rightarrow$Q, b$\rightarrow$R, c$\rightarrow$P, d$\rightarrow$S}
\end{qoptions}
\qfullend
\qn{12} \textit{(Paragraph for Question 12)} For H-atom, the wave function for an electron is $\psi$ = $\dfrac{1}{81}\sqrt{\dfrac{2}{\pi a_0^2}}\cdot\sigma\cdot(6-\sigma)\cdot e^{-\sigma/2}\cdot\cos\theta$, where $\sigma$ = $\dfrac{2\text{Zr}}{\text{na}_0}$ and a$_0$ = 0.529 \AA. Match correctly the parameters given in list-I about the electron and orbital, with its value in list-II.\par
{\footnotesize\setlength{\tabcolsep}{1.5pt}\noindent\begin{tabular}{|c|L{3.7cm}|c|c|}
\hline
 & List-I & & List-II \\
\hline
(a) & Orbital given & (P) & 3s \\
\hline
(b) & Number of radial nodes & (Q) & 1 \\
\hline
(c) & Other orbital with same energy for electron & (R) & 3p$_\text{z}$ \\
\hline
(d) & Distance of radial node from the centre of nucleus & (S) & 4.5a$_0$ \\
\hline
\end{tabular}}\par
\begin{qoptions}
\opt{A}{a-P, b-Q, c-R, d-S}
\opt{B}{a-R, b-Q, c-P, d-S}
\opt{C}{a-R, b-S, c-P, d-Q}
\opt{D}{a-P, b-Q, c-S, d-R}
\end{qoptions}
\qfullbegin
\qn{13} \textit{(Paragraph for Question 13)}\par
\qfig{p2_chem_s3_q12_1}
\qfullsplit
{\footnotesize\setlength{\tabcolsep}{1.5pt}\noindent\begin{tabular}{|L{2.2cm}|c|L{4.3cm}|}
\hline
\textbf{List-I} (Steps in above flow chart) & & \textbf{List-II} (Observation by given step) \\
\hline
(a) \textbf{STEP-1} & (P) & Cu$^{+2}$ is separated from Hg$^{+2}$ \\
\hline
(b) \textbf{STEP-4} & (Q) & Metallic precipitate is obtained \\
\hline
(c) \textbf{STEP-7} & (R) & Dark blue precipitate is formed \\
\hline
(d) \textbf{STEP-8} & (S) & All four metal ions are present in precipitate, No ion in solution \\
\hline
\end{tabular}}\par
The correct code is :-\par
\begin{qoptions}
\opt{A}{a$\rightarrow$S , b$\rightarrow$Q , c$\rightarrow$P , d$\rightarrow$R}
\opt{B}{a$\rightarrow$S , b$\rightarrow$R , c$\rightarrow$P , d$\rightarrow$Q}
\opt{C}{a$\rightarrow$S , b$\rightarrow$P , c$\rightarrow$Q , d$\rightarrow$R}
\opt{D}{a$\rightarrow$P , b$\rightarrow$S , c$\rightarrow$R , d$\rightarrow$Q}
\end{qoptions}
\qfullend
\qn{14} \textit{(Paragraph for Question 14)}\par
\qfig{p2_chem_s3_q13_2}
Given that\par
Atomic mass : P = 31 , Cl = 35.5 , Ag = 108 , S = 32 , O = 16 , H = 1 , Sn = 119\par
{\footnotesize\setlength{\tabcolsep}{1.5pt}\noindent\begin{tabular}{|c|L{2.7cm}|c|c|}
\hline
 & List-I & & List-II \\
\hline
(a) & p is & (P) & 2200 \\
\hline
(b) & (t$_1$ + t$_2$) is & (Q) & 2296 \\
\hline
(c) & r is & (R) & 1044 \\
\hline
(d) & s is & (S) & 3336 \\
\hline
\end{tabular}}\par
The correct code is :-\par
\begin{qoptions}
\opt{A}{a$\rightarrow$S , b$\rightarrow$Q , c$\rightarrow$R , d$\rightarrow$P}
\opt{B}{a$\rightarrow$R , b$\rightarrow$P , c$\rightarrow$Q , d$\rightarrow$S}
\opt{C}{a$\rightarrow$P , b$\rightarrow$S , c$\rightarrow$Q , d$\rightarrow$R}
\opt{D}{a$\rightarrow$S , b$\rightarrow$P , c$\rightarrow$Q , d$\rightarrow$R}
\end{qoptions}
\sect{SECTION-II}
\paraq{1}{2}\par
\qfig{p2_chem_s4_para1_2_1}
Compound (A) was hydrolysed in aqueous acidic medium for some time. After an hour, a mixture of compound (B) and (C) are obtained. properties of (A), (B) and (C) are summarised.\par
{\footnotesize\setlength{\tabcolsep}{1.5pt}\noindent\begin{tabular}{|L{2.6cm}|L{2cm}|c|c|}
\hline
 & (A) & (B) & (C) \\
\hline
Reaction with NaHCO$_3$ & Brisk efferescance & No & No \\
\hline
Blue litmus & Red & No & No \\
\hline
NaOI & No & No & Yellow \\
\hline
{[}Ag(NH$_3$)$_2${]}$^+$ & No & No & Silver mirror \\
\hline
\end{tabular}}\par
\qn{1} 4.4 gm of compound (C) was treated with iodine in aqueous basic medium. Assuming 100\% yield of reaction, calculate mass of tri-iodo compound formed? [At. wt of I = 127, At. wt. of Mg = 24, At wt. of Br = 80]\par
\qn{2} 0.1 mole of compound (B) was reacted with excess of CH$_3$MgBr. Calculate mass of CH$_3$MgBr consumed for complete reaction\par
\qfullbegin
\paraq{3}{4}\par
\qfig{p2_chem_s4_para3_4_1}
\qfullend
\qn{3} Total number of fractions obtained on fractional distillation of all monochloro derivatives product(s) (C)?\par
\qn{4} Calculate mass of Cl$_2$ required to completely react with 9.8 gm of compound (B). Assuming \% yield = 100\%\par
\paraq{5}{6}Cadmium oxide crystallises in NaCl type structure. The compound is however non-stoichiometric with the formula CdO$_{0.95}$. The defect arises due to some cationic positions occupied by neutral Cd atom in place of Cd$^{2+}$ ions and equivalent numbers of anion sites are vacant. (Cd = 112, N$_\text{A}$ = 6 $\times$ 10$^{23}$)\par
\qn{5} What percent of anion sites are vacant ?\par
\qn{6} If the edge length of the unit cell is 400 pm, what would be the density (in gm/cm$^3$) of this non-stoichiometric solid ?\par
\raggedcolumns
\partsec{3}{MATHEMATICS}
\sect{SECTION-I (i)}
\paraq{1}{2}Let f(x) = $\dfrac{1}{1-x}$ and g(x) is defined as g (x) = $\underbrace{\text{f (f (f. . . (f (x)) . . . .))}}_{\text{2n brackets}}$\par
consider\par
\resizebox{\linewidth}{!}{h (x) = $\begin{cases}(\text{g (x)})^6+(\text{g (x)})^3-1 & , \quad x\in\text{R}-\{0,1\}\\ -1 & , \quad x=0\\ 1 & , \quad x=1\end{cases}$}\par
On the basis of above information, answer the following questions :\par
\qn{1} If h(x) is a polynomial then n can be
\begin{qoptions}
\opt{A}{3}
\opt{B}{5}
\opt{C}{6}
\opt{D}{7}
\end{qoptions}
\qn{2} If h(x) is a polynomial then area bounded by y = h(x) and tangent to it at the point where it crosses y-axis is
\begin{qoptions}
\opt{A}{$\dfrac{11}{28}$}
\opt{B}{$\dfrac{3}{28}$}
\opt{C}{$\left|\displaystyle\int_{-1}^{0}\text{h (x) dx}\right|$}
\opt{D}{$\displaystyle\int_{-1}^{0}(-1-\text{h (x)})\,\text{dx}$}
\end{qoptions}
\paraq{3}{4}Let f(x) = $-$x$^2$ + bx + c be a function such that f(x) = f(2 $-$ x) and $\alpha$, $\beta$ be roots of equation f(x) = 0\par
[Note : [k] denotes greatest integer function less than or equal to k]\par
On the basis of above information, answer the following questions :\par
\qn{3} If c $\in$ (0, 1) then minimum value of max ([$\alpha$], [$\beta$]) is equal to
\begin{qoptions}
\opt{A}{2}
\opt{B}{4}
\opt{C}{6}
\opt{D}{8}
\end{qoptions}
\qn{4} If c $\in$ ($-$1,0) then which of the following statement is (are) correct
\begin{qoptions}
\opt{A}{min ([$\alpha$], [$\beta$]) = 0}
\opt{B}{min ([$\alpha$], [$\beta$]) = 2}
\opt{C}{[$\alpha$] + [$\beta$] $<$ 2}
\opt{D}{maximum value of f(x) $\in$ (0, 1)}
\end{qoptions}
\sect{SECTION-I (ii)}
\paraq{5}{6}For all x $\in$ $\left[-\dfrac{7}{2},\dfrac{7}{2}\right]$, let f (x) = $\begin{cases}\|[\text{x}]\| & ;\quad 0\leq\{\text{x}\}<\frac{1}{2}\\ |\text{x}| & ;\quad \frac{1}{2}\leq\{\text{x}\}<1\end{cases}$\par
(where [k] denotes greatest integer function less than or equal to k and \{k\} denotes fractional part of k)\par
On the basis of above information, answer the following questions :\par
\qn{5} Number of solution of the equation f (x) = $\dfrac{5}{2}$ is equal to
\begin{qoptions}
\opt{A}{2}
\opt{B}{3}
\opt{C}{4}
\opt{D}{infinite}
\end{qoptions}
\qn{6} If L denotes number of points of discontinuity and M is the number of points of non-differentiability of function f(x) then (L + M) is equal to
\begin{qoptions}
\opt{A}{18}
\opt{B}{20}
\opt{C}{21}
\opt{D}{22}
\end{qoptions}
\paraq{7}{8}Consider a polynomial for f(x) = a$^2$x$^3$ + ax$^2$ $-$ x + b with local minimum at x = $-$1 with positive extreme values.\par
On the basis of above information, answer the following questions :\par
\qn{7} Let X = \{b : b $\in$ I for a $\in$ I\} and Y = \{b : b $\in$ I for a $\notin$ I\}. Then the maximum value of Z = $\displaystyle\sum_{n\in X}\left(\dfrac{\cos\theta}{2}\right)^n+\sum_{m\in Y}\left(\dfrac{\cos\frac{\theta}{2}}{3}\right)^m$
\begin{qoptions}
\opt{A}{2}
\opt{B}{$\dfrac{7}{2}$}
\opt{C}{$\dfrac{1}{6}$}
\opt{D}{$\dfrac{5}{6}$}
\end{qoptions}
\qn{8} The point of intersection of tangents at (0, f(0)) and (f(0), f(f(0))) for f(x) = a$^2$x$^3$ + ax$^2$ $-$ x + n where a $\in$ I and n is the corresponding smallest positive integral value of b is
\begin{qoptions}
\opt{A}{$\left(\dfrac{9}{5},-\dfrac{4}{5}\right)$}
\opt{B}{(2, 0)}
\opt{C}{$\left(\dfrac{3}{5},\dfrac{2}{5}\right)$}
\opt{D}{(2, $-$1)}
\end{qoptions}
\paraq{9}{10}Consider the parabola x$^2$ = 4y and the circle x$^2$ + y$^2$ $-$ 10y + 25 = r$^2$(r $>$ 0). Given that the circle touches the parabola at the points P and Q. Let R be the point of intersection of tangents to parabola at P and Q and S be the centre of circle.\par
On the basis of above information, answer the following questions :\par
\qn{9} Which of the following statement(s) is/are TRUE ?
\begin{qoptions}
\opt{A}{Centroid of $\Delta$PQR is $\left(0,\dfrac{7}{9}\right)$}
\opt{B}{Distance between chord PQ and directrix of parabola is 4}
\opt{C}{Area of $\Delta$PQR is $\dfrac{40}{3\sqrt{3}}$}
\opt{D}{Common tangents to circle and parabola intersect at (0, 3)}
\end{qoptions}
\qn{10} The equation of circle which passes through the vertex of the parabola x$^2$ = 4y and touches it at the point M($-$4, 4) is
\begin{qoptions}
\opt{A}{x$^2$ + y$^2$ $-$ 10x $-$ 18y = 0}
\opt{B}{x$^2$ + y$^2$ $-$ 8x $-$ 16y = 0}
\opt{C}{x$^2$ + y$^2$ + 8x = 0}
\opt{D}{x$^2$ + y$^2$ $-$ 8y = 0}
\end{qoptions}
\sect{SECTION-I (iii)}
\paraq{11}{11}Let f (x) = $\displaystyle\lim_{n\to\infty}n\left(\text{x}^{\frac{2022}{n}}-\text{x}^{\frac{2023}{n}}\right)$ ; x $>$ 0.\par
\qn{11} Match List-I with List-II and select the correct answer using the code given below the list.\par
{\footnotesize\setlength{\tabcolsep}{1.5pt}\noindent\begin{tabular}{|c|L{3.6cm}|c|c|}
\hline
 & List-I & & List-II \\
\hline
(a) & Number of solution of equation f(|x|) = e$^{-2023x^2}$ & (P) & 1 \\
\hline
(b) & Number of solution of equation f(|x|) = e$^{2023x^2}$ & (Q) & 2 \\
\hline
(c) & Number of solution of equation f(x) = e$^{2022}$ & (R) & 3 \\
\hline
(d) & Number of solution of equation f(x) = e$^{-2022}$ & (S) & 0 \\
\hline
\end{tabular}}\par
\begin{qoptions}
\opt{A}{a $\rightarrow$ Q; b $\rightarrow$ Q; c $\rightarrow$ P; d $\rightarrow$ Q}
\opt{B}{a $\rightarrow$ P; b $\rightarrow$ P; c $\rightarrow$ Q; d $\rightarrow$ S}
\opt{C}{a $\rightarrow$ P; b $\rightarrow$ Q; c $\rightarrow$ P; d $\rightarrow$ Q}
\opt{D}{a $\rightarrow$ Q; b $\rightarrow$ P; c $\rightarrow$ Q; d $\rightarrow$ S}
\end{qoptions}
\paraq{12}{12}Consider a 3 $\times$ 3 matrix A whose element\par
a$_\text{ij}$ = tan$^{-1}$(tan(i + j)), 1 $\leq$ i, j $\leq$ 3\par
Tr(X) denotes trace of matrix X\par
$|$X$|$ denotes determinant of X\par
adj(X) denotes adjoint of matrix X, [.] = greatest integer function\par
If $|$A$|$ = $\ell\pi^\alpha$ $-$ m$\pi^\beta$ $-$ n$\pi^\gamma$ ($\alpha>\beta>\gamma$)\par
\qn{12} Match List-I with List-II and select the correct answer using the code given below the list.\par
{\footnotesize\setlength{\tabcolsep}{1.5pt}\noindent\begin{tabular}{|c|c|c|c|}
\hline
 & List-I & & List-II \\
\hline
(a) & [Tr(A)] & (P) & $-$1 \\
\hline
(b) & $\alpha+\beta+\gamma$ & (Q) & 6 \\
\hline
(c) & $\ell$ + m + n & (R) & 0 \\
\hline
(d) & [a$_{13}$] & (S) & 1 \\
\hline
\end{tabular}}\par
\begin{qoptions}
\opt{A}{a $\rightarrow$ Q; b $\rightarrow$ P; c $\rightarrow$ R; d $\rightarrow$ S}
\opt{B}{a $\rightarrow$ P; b $\rightarrow$ Q; c $\rightarrow$ Q; d $\rightarrow$ R}
\opt{C}{a $\rightarrow$ P; b $\rightarrow$ R; c $\rightarrow$ Q; d $\rightarrow$ S}
\opt{D}{a $\rightarrow$ S; b $\rightarrow$ P; c $\rightarrow$ Q; d $\rightarrow$ R}
\end{qoptions}
\paraq{13}{13}Let L = $\displaystyle\lim_{n\to\infty}\dfrac{n^3\left(e^{\frac{1}{n}}+e^{\frac{2}{n}}+\ldots+e\right)}{(n+1)^m\cdot\left(1^{2m}+2^{2m}+\ldots n^{2m}\right)}$ and\par
\resizebox{\linewidth}{!}{M = $\displaystyle\lim_{x\to0}\dfrac{\sqrt{1+x^2}\tan\left(\sin\left(\tan^{-1}x\right)\right)+2\sqrt{1-x^2}\sin\left(\tan\left(\sin^{-1}x\right)\right)-3x}{x^P}$}\par
are non zero finite real numbers (where [.] greatest integer function)\par
\qn{13} Match List-I with List-II and select the correct answer using the code given below the list.\par
{\footnotesize\setlength{\tabcolsep}{1.5pt}\noindent\begin{tabular}{|c|c|c|c|}
\hline
 & List-I & & List-II \\
\hline
(a) & m & (P) & 1 \\
\hline
(b) & [L] & (Q) & 3 \\
\hline
(c) & P & (R) & 5 \\
\hline
(d) & [M] & (S) & $-$1 \\
\hline
\end{tabular}}\par
\begin{qoptions}
\opt{A}{a $\rightarrow$ P; b $\rightarrow$ R; c $\rightarrow$ R; d $\rightarrow$ S}
\opt{B}{a $\rightarrow$ S; b $\rightarrow$ P; c $\rightarrow$ R; d $\rightarrow$ Q}
\opt{C}{a $\rightarrow$ S; b $\rightarrow$ R; c $\rightarrow$ P; d $\rightarrow$ R}
\opt{D}{a $\rightarrow$ R; b $\rightarrow$ S; c $\rightarrow$ P; d $\rightarrow$ Q}
\end{qoptions}
\paraq{14}{14}If I = $\displaystyle\int_0^\infty e^{-x^2}dx$ = $\alpha$ and $\displaystyle\int_0^\infty\dfrac{\sin x}{x}dx$ = $\beta$\par
\qn{14} Match List-I with List-II and select the correct answer using the code given below the list.\par
{\footnotesize\setlength{\tabcolsep}{1.5pt}\noindent\begin{tabular}{|c|L{3.9cm}|c|c|}
\hline
 & List-I & & List-II \\
\hline
(a) & $\displaystyle\int_0^\infty\dfrac{\sin ax\cos bx}{x}dx$ (a, b $>$ 0) & (P) & $\alpha$ \\
\hline
(b) & $\displaystyle\int_0^\infty\dfrac{e^{-x}}{\sqrt{x}}dx$ & (Q) & 2$\alpha$ \\
\hline
(c) & $\displaystyle\int_0^\infty x^2e^{-x^2}dx$ & (R) & $\beta$ \\
\hline
(d) & $\displaystyle\int_0^\infty\dfrac{\sin^2x}{x^2}dx$ & (S) & $\dfrac{\alpha}{2}$ \\
\hline
\end{tabular}}\par
\begin{qoptions}
\opt{A}{a $\rightarrow$ R; b $\rightarrow$ Q; c $\rightarrow$ S; d $\rightarrow$ R}
\opt{B}{a $\rightarrow$ R; b $\rightarrow$ R; c $\rightarrow$ Q; d $\rightarrow$ S}
\opt{C}{a $\rightarrow$ S; b $\rightarrow$ Q; c $\rightarrow$ R; d $\rightarrow$ R}
\opt{D}{a $\rightarrow$ P; b $\rightarrow$ Q; c $\rightarrow$ R; d $\rightarrow$ S}
\end{qoptions}
\sect{SECTION-II}
\paraq{1}{2}Suppose A, B, C are three points corresponding to complex numbers z$_0$ = ai, z$_1$ = $\dfrac{1}{2}$ + bi, z$_2$ = 1 + ci (a, b, c $\in$ R) respectively. Now consider the curve z = z$_0$cos$^4$t + 2z$_1$cos$^2$t sin$^2$t + z$_2$ sin$^4$t (t $\in$ R)\par
$\omega_1$ represents z when A, B, C are non-collinear and $\omega_2$ represents z when A, B, C are collinear.\par
\qn{1} If b $<$ a (a, b $\in$ R$^+$) the number of point of intersections of $\omega_2$ with the curve z = t + ibt$^2$ (t $\in$ R)\par
\qn{2} If b $>$ a (a, b $\in$ R$^+$) the number of point of intersections of $\omega_2$ with the curve z = t + ibt$^2$ (t $\in$ R)\par
\paraq{3}{4}Let f : R $\rightarrow$ R be a function defined as f(x) = 3x$^5$ $-$ 25x$^3$ + 60x + 5 and\par
g (x) = $\begin{cases}\max\{\text{f (t)};-4\leq\text{t}\leq\text{x}\} & ;\quad -4\leq\text{x}\leq0\\ \min\{\text{f (t)};0<\text{t}\leq\text{x}\} & ;\quad 0<\text{x}\leq2\\ \text{f (x)}-16 & ;\quad \text{x}>2\end{cases}$\par
\qn{3} Total number of points at which the graph of y = f(x) and y = g(x) is non-differentiable in [$-$4, $\infty$) is/are\par
\qn{4} If the equation f(x) + $\lambda$ = 0 is having exactly three distinct real roots then total number of possible integral values of $\lambda$ are\par
\paraq{5}{6}There exists a matrix Q such that PQP$^\text{T}$ = N where P = $\begin{bmatrix}1&2&0\\2&1&0\\0&0&1\end{bmatrix}$. Given N is a diagonal matrix of form\par
N = diag(n$_1$, n$_2$, n$_3$) where n$_1$, n$_2$, n$_3$ are three values of n satisfying equation det(P $-$nI) = 0, n$_1$ $<$ n$_2$ $<$ n$_3$\par
(I is identity matrix of order 3)\par
\qn{5} The trace of matrix Q$^2$ is equal to\par
\qn{6} The trace of matrix P$^{2022}$ is 3$^{2022}$ + k then $\dfrac{\text{k}}{4}$ is\par
\end{multicols}
\end{document}
